\subsection{Persistenz mit PostgreSQL, SQLAlchemy und Alembic}
\label{subsec:persistenz-postgresql}

PostgreSQL ist keine Pflichtkomponente und kein eigenes Profil. Die wählbare
Capability \texttt{postgres} setzt \texttt{database} voraus, das wiederum ein
Backend benötigt; Web-only und Desktop-local werden deshalb abgewiesen. Die
generische Datenbankfähigkeit enthält SQLAlchemy-Engine, Sessions und Alembic,
während die PostgreSQL-Erweiterung den Psycopg-3-Treiber, URL-Vorgaben und
Integrationstests ergänzt. PostgreSQL selbst bleibt eine getrennte
Laufzeiteinheit \parencite{postgresql2026docs,kleiveist2026database}.

\path{backend/app/db/engine.py} erzeugt die Engine erst beim ersten Zugriff und
aktiviert \texttt{pool\_pre\_ping}; Sessions werden über eine Generator-Dependency
geschlossen. Die leere deklarative \texttt{Base} ist nur ein Erweiterungspunkt
für spätere Fachmodelle. Das entspricht der von SQLAlchemy dokumentierten
Trennung von Engine, Pool, Dialekt und Session sowie dem verzögerten Aufbau der
ersten Verbindung \parencite{sqlalchemy2026engines}.

Alembic verwendet dieselbe Metadatenbasis für Online- und Offline-Migrationen.
Da das Template keine Fachmodelle enthält, existiert noch keine initiale
Revision; weder \texttt{create\_all()} noch automatische Migrationen beim
FastAPI-Start sind implementiert. Schemaänderungen erfolgen ausschließlich über
explizite Alembic-Kommandos, deren Umgebung und Revisionsablauf die offizielle
Dokumentation beschreibt \parencite{alembic2026tutorial,kleiveist2026database}.
\texttt{/api/health} bleibt datenbankunabhängig, während
\texttt{/api/ready} bei aktivierter Datenbankfähigkeit die Erreichbarkeit prüft.
